%==============================================================
%  TÓM TẮT / ABSTRACT
%==============================================================

\chapter*{Tóm tắt}
\addcontentsline{toc}{chapter}{Tóm tắt}
\markboth{Tóm tắt}{Tóm tắt}

Báo cáo trình bày một cách hệ thống cơ sở toán học, thuật toán và thực nghiệm
của \emph{mạng nơ-ron thông tin vật lý} (physics-informed neural network, PINN)
khi được dùng làm công cụ giải phương trình vi phân.

Điểm khác biệt về trọng tâm so với tài liệu học sâu tổng quát là báo cáo khảo
sát mạng truyền thẳng từ góc độ \emph{đạo hàm theo biến đầu vào} chứ không phải
theo tham số. Hai hệ thức truy hồi cho Jacobi và Hessian
(Mệnh đề~\ref{md:jacobian}, \ref{md:hessian}) được thiết lập tường minh và trở
thành công cụ trung tâm cho mọi tính toán phần dư về sau; từ chúng suy ra ngay
một tiêu chuẩn \emph{cần} đối với hàm kích hoạt. Trường hợp ReLU được phát biểu
chính xác ở dạng mạnh hơn nhiều một sự ``suy giảm'': hàm mục tiêu phần dư của
phương trình bậc hai là \emph{hằng số địa phương}, nên mọi thuật toán dựa trên
gradient đứng yên tuyệt đối (Hệ quả~\ref{hq:pinn-hong}).

Về phương pháp, báo cáo trả lời câu hỏi thường bị bỏ qua: \emph{vì sao cực tiểu
hoá phần dư lại là việc đáng làm}. Câu trả lời là một chặn ổn định của bài toán
vi phân (Mệnh đề~\ref{md:chan-on-dinh}), và hằng số của chặn ấy suy biến theo
$1/\nu$ với bài toán nhiễu kỳ dị (Mệnh đề~\ref{md:chan-suy-bien}). Kết hợp với
bệnh lý gradient ở tầng tối ưu và thiên kiến phổ ở tầng xấp xỉ, báo cáo tổng
hợp \emph{ba nguyên nhân độc lập, nằm ở ba tầng khác nhau}, cùng làm PINNs suy
giảm trên đúng lớp bài toán có thang độ dài nhỏ (Bảng~\ref{tab:ba-nguyen-nhan}).

Phần thực nghiệm cài đặt PINNs bằng PyTorch, không dùng thư viện đóng gói sẵn,
và kiểm chứng từng phát biểu lý thuyết bằng mười một thí nghiệm có tiêu chí thành bại
xác định trước. Kết quả được đối chiếu với lý thuyết ở
Bảng~\ref{tab:doi-chieu}. Những kết quả không đạt kỳ vọng cũng được trình bày
đầy đủ: \emph{năm} kết quả buộc phải siết lại chính các chương lý thuyết --- một
phát biểu về thiên kiến phổ (Nhận xét~\ref{nx:hieu-chinh}), một cách đọc hằng số
ổn định như thể nó là cơ chế đang vận hành (Nhận xét~\ref{nx:doc-tn7}), một
ước lượng chi phí đồ thị thấp hơn giá trị đo được
(Nhận xét~\ref{nx:chi-phi-bac-cao}), và hai phát biểu về bài toán ngược
(Nhận xét~\ref{nx:doc-tn11}) --- còn một kết quả tiêu cực về thuật toán
cân bằng trọng số thì dẫn tới một chứng minh mới về chế độ phản hồi dương của
quy tắc ấy (Bổ đề~\ref{bd:phan-hoi-duong}). Việc lặp lại mọi thí nghiệm có phụ thuộc ngẫu
nhiên trên năm hạt giống còn loại bỏ được bốn con số mà chính báo cáo đã trích
dẫn trước đó, trong khi các kết luận định tính đều đứng vững
(Mục~\ref{sec:tn8}). Phép đo chi phí cho một kết quả cần được nêu
thẳng: trên bài toán một chiều được khảo sát, PINN chậm hơn bộ giải sai phân hữu
hạn khoảng bốn bậc độ lớn ở cùng mức sai số, nên giá trị của phương pháp phải được biện minh
bằng bài toán ngược, dữ liệu biên khuyết và số chiều cao chứ không bằng chi phí
(Mục~\ref{sec:tn9}). Cuối cùng, áp dụng chính các biện pháp khắc phục mà báo cáo
đề xuất --- cấu hình của công trình gốc cùng lấy mẫu thích nghi --- đưa sai số
Burgers trung vị từ $1{,}38\times10^{-1}$ xuống $6{,}48\times10^{-4}$, tức tốt hơn
$213$ lần và cùng bậc với công trình gốc, nhưng khoảng cách về chi phí vẫn ở bậc
$10^{4}$ (Mục~\ref{sec:tn10}). Một thí nghiệm tách biến cho thấy mức cải thiện ấy
đến gần như hoàn toàn từ cấu hình mạng và tối ưu hoá: ba chiến lược lấy mẫu
thích nghi được thử đều chỉ phân phối lại sai số giữa lớp sốc và vùng trơn, chứ
không làm giảm tổng sai số. Thí nghiệm duy nhất đặt PINN vào đúng động lực của nó
--- tìm độ nhớt từ $2\,000$ điểm đo có nhiễu, không biết điều kiện đầu và biên ---
cho sai số trung vị dưới $0{,}5\%$, tái lập kết quả của công trình gốc; nhưng một
phương pháp liên hợp rời rạc được cho cùng thông tin đạt độ chính xác cùng bậc và
có đáp số sớm hơn, nên ưu thế của PINN ở đây nằm ở công sức xây dựng phương pháp
chứ không ở chi phí tính toán (Mục~\ref{sec:tn11}).

\bigskip
\noindent\textbf{Từ khoá:} mạng nơ-ron thông tin vật lý, phương trình vi phân,
vi phân tự động, xấp xỉ phổ quát, thiên kiến phổ, bệnh lý gradient, chặn ổn
định, L-BFGS.

%--------------------------------------------------------------
\chapter*{Abstract}
\addcontentsline{toc}{chapter}{Abstract}
\markboth{Abstract}{Abstract}

\begin{english}
This report gives a systematic account of the mathematical, algorithmic and
experimental foundations of \emph{physics-informed neural networks} (PINNs) used
as solvers for differential equations.

Its focus differs from that of general deep-learning texts: feedforward networks
are studied through their derivatives \emph{with respect to the input} rather
than to the parameters. Two explicit recurrences for the input Jacobian and
Hessian are established and used throughout for residual computations; from them
follows a \emph{necessary} condition on the activation function. The ReLU case is
stated in a form considerably stronger than mere degradation: for a second-order
operator the residual objective is \emph{locally constant}, so every
gradient-based algorithm is exactly stationary.

Methodologically, the report answers a question that is usually skipped: \emph{why
minimising the residual is worth doing at all}. The answer is a stability bound
for the underlying differential problem, whose constant degenerates as $1/\nu$ for
singularly perturbed problems. Combined with gradient pathology at the
optimisation level and spectral bias at the approximation level, this yields
\emph{three independent causes, situated at three different levels}, that jointly
degrade PINNs on precisely the class of problems with small length scales.

The experimental part implements PINNs in PyTorch without any packaged PINN
library and tests each theoretical claim through eleven experiments with success
criteria fixed in advance. Negative results are reported in full: five of them
forced a restatement in the theoretical chapters themselves -- a claim about
spectral bias, a reading of the stability constant as an operative mechanism,
a cost estimate that measurement showed to be too low, and two claims about
inverse problems -- and a negative
result on adaptive loss balancing led to a new proof that the balancing rule has
a positive-feedback failure mode. Repeating every seed-dependent experiment over
five random seeds further eliminated four figures the report had itself quoted,
while leaving all of its qualitative conclusions intact. A
direct cost measurement gives a further result worth stating plainly: on the
one-dimensional problems studied here the PINN is about four orders of magnitude
slower than a finite-difference solver at the same accuracy, so the value of the
method must be argued from inverse problems, incomplete boundary data and high
dimension rather than from cost. Finally, applying the remedies the report itself proposes -- the
original authors' configuration together with residual-based adaptive
refinement -- lowers the median Burgers error from $1.38\times10^{-1}$ to
$6.48\times10^{-4}$, a $213$-fold improvement on par with the original work,
while the cost gap stays at four orders of magnitude. An ablation shows the gain
comes almost entirely from the network and optimiser configuration: the three
adaptive sampling strategies tested only redistribute the error between the
shock and the smooth region, without reducing it overall. The one experiment that
places the PINN in its intended setting -- identifying the viscosity from $2\,000$
noisy measurements without initial or boundary data -- gives a median error below
$0.5\%$, reproducing the original work; but a discrete adjoint method given the
same information reaches comparable accuracy and settles on its answer sooner, so
the PINN's advantage there lies in the effort of building the method rather than
in computational cost.

\bigskip
\noindent\textbf{Keywords:} physics-informed neural networks, differential
equations, automatic differentiation, universal approximation, spectral bias,
gradient pathology, stability bound, L-BFGS.
\end{english}
